\section{First-Order DDP Derivation}
\label{app:user_ddp}

The copper-plate problem below is included as the reference problem used to
derive and test the first-order Differential Dynamic Programming workflow. It
keeps the temporal battery coupling explicit without introducing network
equations, so the information passed between period subproblems can be seen
directly.

\section{Problem Formulations}
\label{sec:copper_plate}

\subsection{Copper-Plate MPOPF}

A storage device and grid connection serve demand $p_L^t$ over
$\mathcal T=\{1,\ldots,T\}$.  Here $P_B^t>0$ discharges, and $B^t$ is the
end-of-period energy; hence $B^0$ is initial and $B^T$ is terminal.  The sign
convention follows~\cite{Jha2025}.
\noindent\fbox{\begin{minipage}{\dimexpr\columnwidth-2\fboxsep-2\fboxrule\relax}
\small
\begin{subequations}
\label{eq:cp_all}
\begin{alignat}{2}
\min_{P_{\text{Subs}},\,P_B,\,B} \;\;
  & \sum_{t=1}^{T} \Bigl[ c^t P_{\text{Subs}}^t
                        + C_B \bigl(P_B^t\bigr)^2 \Bigr] \Delta t \nonumber \\[-2pt]
  & \hspace{34pt} +\; w \bigl(B^{T} - B_{\text{init}}\bigr)^2
\label{eq:cp_obj} \\[2pt]
\text{s.t.} \;\;
  & P_{\text{Subs}}^t + P_B^t - p_L^t = 0, &\quad& t \in \mathcal{T}
\label{eq:cp_balance} \\
  & B^0 = B_{\text{init}}
\label{eq:cp_soc_init} \\
  & B^{t} = B^{t-1} - P_B^t \Delta t, && t \in \mathcal{T}
\label{eq:cp_soc_dyn} \\
  & P_{\text{Subs}}^t \geq 0, && t \in \mathcal{T}
\label{eq:cp_pg_lim} \\
  & \underline{P_B} \leq P_B^t \leq \overline{P_B}, && t \in \mathcal{T}
\label{eq:cp_pb_lim} \\
  & \underline B \leq B^{t} \leq \overline B, && t \in \mathcal{T}
\label{eq:cp_soc_lim}
\end{alignat}
\end{subequations}
\end{minipage}}
The \emph{base instance} is
\eqref{eq:cp_obj}--\eqref{eq:cp_soc_dyn}, with all three bounds dropped.

\section{SOCP MPOPF Formulation}
\label{sec:network_mpopf}

Let $\mathcal N$, $\mathcal L$, $\mathcal B$, and $\mathcal D$ denote buses,
directed branches, battery buses, and PV buses; $j_{\mathrm{subs}}=1$ is the
substation bus and $\mathcal C(j)$ contains the children of $j$. The network
MPOPF is

\noindent\fbox{\begin{minipage}{\dimexpr\columnwidth-2\fboxsep-2\fboxrule\relax}
\footnotesize
\setlength{\jot}{1pt}
\begin{align}
\min\quad
&\sum_{t\in\mathcal T} c^t S_{\mathrm{base}}P_{\mathrm{Subs}}^t\Delta t
\nonumber\\[-1pt]
&+\sum_{t\in\mathcal T}\sum_{j\in\mathcal B}
 C_B\bigl(S_{\mathrm{base}}P_{B,j}^t\bigr)^2\Delta t
\nonumber\\[-1pt]
&+w\sum_{j\in\mathcal B}\bigl(B_j^T-B_{\mathrm{init},j}\bigr)^2
\label{eq:net_obj}\\
\mathrm{s.t.}\quad
&P_{\mathrm{Subs}}^t-\sum_{(j_{\mathrm{subs}},k)\in\mathcal L}
 P_{j_{\mathrm{subs}}k}^t=0,
\label{eq:net_sub_p}\\
&P_{\mathrm{Subs}}^t\geq0,\\
&Q_{\mathrm{Subs}}^t-\sum_{(j_{\mathrm{subs}},k)\in\mathcal L}
 Q_{j_{\mathrm{subs}}k}^t=0,
\label{eq:net_sub_q}\\
&\sum_{k\in\mathcal C(j)}P_{jk}^t-P_{ij}^t+r_{ij}\ell_{ij}^t
 =P_{B,j}^t+p_{D,j}^t-p_{L,j}^t,
\label{eq:net_p_balance}\\
&\sum_{k\in\mathcal C(j)}Q_{jk}^t-Q_{ij}^t+x_{ij}\ell_{ij}^t
 =q_{D,j}^t-q_{L,j}^t,
\label{eq:net_q_balance}\\
&v_j^t=v_i^t-2(r_{ij}P_{ij}^t+x_{ij}Q_{ij}^t)
 +(r_{ij}^2+x_{ij}^2)\ell_{ij}^t,
\label{eq:net_voltage}\\
&(P_{ij}^t)^2+(Q_{ij}^t)^2\leq v_i^t\ell_{ij}^t,
\label{eq:net_soc}\\
&\ell_{ij}^t\geq0,\\
&v_{\mathrm{Subs}}^t=1.05^2,\\
&\underline V_j^2\leq v_j^t\leq\overline V_j^2,
\label{eq:net_v_bounds}\\
&-q_{D,j}^{\max,t}\leq q_{D,j}^t\leq q_{D,j}^{\max,t},
\label{eq:net_pv_bounds}\\
&q_{D,j}^{\max,t}=\sqrt{S_{D,j}^2-(p_{D,j}^t)^2},\\
&B_j^t=B_j^{t-1}-P_{B,j}^t\Delta t,
\label{eq:net_b_dyn}\\
&B_j^0=B_{\mathrm{init},j},\\
&\mathrm{soc}_j^{\min}B_{R,j}\leq B_j^t
 \leq\mathrm{soc}_j^{\max}B_{R,j},
\label{eq:net_b_bounds}\\
&-P_{B,j}^{R}\leq P_{B,j}^t\leq P_{B,j}^{R}.
\label{eq:net_pb_bounds}
\end{align}
\end{minipage}}


\section{Simulation Setup}

\begin{table}[H]
\centering
\caption{Copper-Plate Instance Data}
\label{tab:cp_data}
\setlength{\tabcolsep}{5pt}
\begin{tabular}{lcccccc}
\toprule
\textbf{Interval $t$} & $1$ & $2$ & $3$ & $4$ & $5$ & $6$ \\
\midrule
Demand $p_L^t$     & $1.657$ & $1.888$ & $1.998$ & $1.834$ & $1.623$ & $1.657$ \\
Energy price $c^t$ & $0.140$ & $0.197$ & $0.175$ & $0.105$ & $0.083$ & $0.140$ \\
\bottomrule
\end{tabular}

\smallskip

\renewcommand{\arraystretch}{1.15}
\begin{tabular}{lc@{\hskip 10pt}lc}
\toprule
Period length $\Delta t$ & $1$   & Initial energy $B_{\text{init}}$ & $2.0$ \\
Throughput $C_B$         & $0.05$ & Terminal weight $w$  & $5.0$ \\
\bottomrule
\end{tabular}

\smallskip

\renewcommand{\arraystretch}{1.25}
\begin{tabular}{lc}
\toprule
\textit{Bounds, constrained variants only} & \\
\midrule
$\underline{P_B},\,\overline{P_B}$ & $-0.45,\;0.45$ \\
$\underline B,\,\overline B$       & $1.20,\;1.95$ \\
\bottomrule
\end{tabular}

\smallskip

\raggedright
\footnotesize All quantities are in per unit on a 1000-kVA base, $1$~p.u.\ $=
1000$~kW and $1$~p.u.h $= 1000$~kWh. Bounds apply only to the constrained
variants.
\end{table}


The derivation is stated directly in MPOPF notation. Equality multipliers are denoted by $\lambda$ and
inequality multipliers by $\mu$.  Thus the dynamics multiplier called $\mu[t]$
in the earlier derivation is written $\lambda_{\mathrm{soc}}^t$ here.

\subsection{Centralized Lagrangian}

\noindent\fbox{\begin{minipage}{\dimexpr\columnwidth-2\fboxsep-2\fboxrule\relax}
Write every inequality as $g\leq0$.  The Lagrangian of
\eqref{eq:cp_all} is
\begin{align}
\mathcal L={}&\sum_{t=1}^{T}\!\left[c^tP_{\mathrm{Subs}}^t
 +C_B(P_B^t)^2\right]\Delta t
 +w(B^T-B_{\mathrm{init}})^2 \nonumber\\
&+\sum_{t=1}^{T}\lambda_{\mathrm{bal}}^t
 (P_{\mathrm{Subs}}^t+P_B^t-p_L^t) \nonumber\\
&+\sum_{t=1}^{T}\lambda_{\mathrm{soc}}^t
 (B^t-B^{t-1}+\Delta t P_B^t) \nonumber\\
&+\sum_{t=1}^{T}\mu_{\underline{P_{\mathrm{Subs}}}}^t
 (-P_{\mathrm{Subs}}^t) \nonumber\\
&+\sum_{t=1}^{T}\!\left[\mu_{\underline{P_B}}^t(\underline{P_B}-P_B^t)
 +\mu_{\overline{P_B}}^t(P_B^t-\overline{P_B})\right] \nonumber\\
&+\sum_{t=1}^{T}\!\left[\mu_{\underline B}^t(\underline B-B^t)
 +\mu_{\overline B}^t(B^t-\overline B)\right].
\label{eq:app_lagrangian}
\end{align}
\end{minipage}}

\smallskip
\noindent\fbox{\begin{minipage}{\dimexpr\columnwidth-2\fboxsep-2\fboxrule\relax}
\noindent\textbf{Inter-period optimality condition.}\quad
Here $B^0=B_{\mathrm{init}}$ is fixed. The local SPOPF solver enforces primal
and dual feasibility, complementary slackness, and stationarity with respect
to its own power variables. The condition that couples adjacent periods is
stationarity with respect to $B^t$, for $t=1,\ldots,T-1$:
\begin{align}
0={}&\lambda_{\mathrm{soc}}^t-\lambda_{\mathrm{soc}}^{t+1}
 -\mu_{\underline B}^t+\mu_{\overline B}^t,
\label{eq:app_kkt_b}
\end{align}
Equation~\eqref{eq:app_kkt_b} shows why the successor multiplier must enter the
preceding period. At $t=T$, the terminal penalty already included in the final
SPOPF terminates this chain.
\end{minipage}}

\newpage
\vspace*{5.5\baselineskip}
\subsection{Per-Period Decomposition}

\noindent\fbox{\begin{minipage}{\dimexpr\columnwidth-2\fboxsep-2\fboxrule\relax}
At sweep $k$, period $t$ receives $\textcolor{cGiven}{\boldsymbol{B^{k,t-1}}}$ from the period just solved and
the successor multiplier $\textcolor{cFuture}{\boldsymbol{\lambda_{\mathrm{soc}}^{k-1,t+1}}}$ from the previous
sweep.  It solves
\begin{align}
\min_{P_{\mathrm{Subs}}^t,P_B^t,B^t}\quad
&c^tP_{\mathrm{Subs}}^t\Delta t+C_B(P_B^t)^2\Delta t
-\textcolor{cFuture}{\boldsymbol{\lambda_{\mathrm{soc}}^{k-1,t+1}}}B^t
\label{eq:app_stage_obj}\\
\mathrm{s.t.}\quad
&P_{\mathrm{Subs}}^t+P_B^t=p_L^t,
\label{eq:app_stage_bal}\\
&B^t-\textcolor{cGiven}{\boldsymbol{B^{k,t-1}}}+\Delta tP_B^t=0,
\label{eq:app_stage_soc}\\
&P_{\mathrm{Subs}}^t\geq0,\quad
\underline{P_B}\leq P_B^t\leq\overline{P_B}, \nonumber\\
&\underline B\leq B^t\leq\overline B.
\label{eq:app_stage_limits}
\end{align}
For $t=T$, the last term of \eqref{eq:app_stage_obj} is replaced by
$w(B^T-B_{\mathrm{init}})^2$.  After solving, the equality multiplier is read
from the local SOC equation,
\begin{equation}
\lambda_{\mathrm{soc}}^{k,t}
=-\operatorname{dual}\!\left(B^t-\textcolor{cGiven}{\boldsymbol{B^{k,t-1}}}+\Delta tP_B^t=0\right).
\label{eq:app_dual_read}
\end{equation}
The term $-\textcolor{cFuture}{\boldsymbol{\lambda_{\mathrm{soc}}^{k-1,t+1}}}B^t$ inserts the successor term
missing from a standalone OPF's stationarity in $B^t$.  The earlier program
wrote the equivalent expression
\begin{equation}
\textcolor{cFuture}{\boldsymbol{\lambda_{\mathrm{soc}}^{k-1,t+1}}}
\left(B^{k-1,t+1}-B^t+\Delta tP_B^{t+1}\right).
\label{eq:app_original_coupling}
\end{equation}
Because $B^{k-1,t+1}$ is fixed and the auxiliary $P_B^{t+1}$ is not the next
period's actual decision, only the $-\textcolor{cFuture}{\boldsymbol{\lambda_{\mathrm{soc}}^{k-1,t+1}}}B^t$
part supplies the intended inter-period sensitivity.
\end{minipage}}

\subsection{Sweep}

The workflow is
\begin{equation}
t=1\rightarrow T:\quad
\textcolor{cGiven}{\boldsymbol{B^{k,t-1}}}\rightarrow\eqref{eq:app_stage_obj}
\rightarrow(B^{k,t},\lambda_{\mathrm{soc}}^{k,t}),
\qquad k\leftarrow k+1,
\label{eq:app_sweep}
\end{equation}
and stops when
\begin{equation}
\max\!\left(\lVert B^k-B^{k-1}\rVert_2,
\lVert\lambda_{\mathrm{soc}}^k-\lambda_{\mathrm{soc}}^{k-1}\rVert_2\right)
<\epsilon.
\label{eq:app_stop}
\end{equation}
This passes only the first derivative of the future cost.  FilterDDP also
passes its curvature $V_{BB}^t$; the first-order workflow has no corresponding
quantity.

% ---------------------------------------------------------------------------
% What the first-order DDP scheme exchanges, and how slowly. Drawn for T = 4
% stages over 4 sweeps -- the smallest grid that shows terminal information
% completing its walk from the last stage to the first.
%
% Reading the layout: ROWS are sweeps (one pass of the outer loop), COLUMNS are
% stages. One cell is one small QP. The two couplings are deliberately drawn at
% different angles because that difference IS the mechanism:
%   horizontal, solid   B[t-1] from THIS sweep -- a hard constraint, so it moves
%                       left to right within a row and is always current
%   diagonal, dashed    mu[t+1] and B[t+1] from the PREVIOUS sweep -- a soft
%                       linear term, so it can only step one column left per row
% The diagonal is why the terminal stage's information needs T sweeps to reach
% stage 1; the highlighted staircase traces exactly that walk.
%
% Colour semantics follow the tADMM exchange figure of the companion study, so
% the same narrative carries across both:
%   blue  = local / primal, owned by the sweep that solves it   (B[t])
%   red   = stale, owned by no sweep in progress                (B[t+1] frozen)
%   green = dual, the feedback that closes the loop             (mu[t+1])
%
% GREYSCALE PROOFING. Hue is never the only cue:
%   solved cell     solid dark fill, sharp corners, heavy border
%   stale parameter near-white fill, DASHED border
%   dual            DASHED line only, never a fill, so it cannot collide with
%                   red under deuteranopia
%   the staircase   drawn heavier than every other arrow
% ---------------------------------------------------------------------------
\definecolor{ddpLoc}{HTML}{1F5186}     % deep blue -- primal, solved this sweep
\definecolor{ddpLocE}{HTML}{143A61}
\definecolor{ddpStale}{HTML}{C0392B}   % red       -- frozen from last sweep
\definecolor{ddpStaleE}{HTML}{8E2A20}
\definecolor{ddpDual}{HTML}{0F6B45}    % green     -- the dual mu

\begin{figure*}[!t]
\centering
\begin{tikzpicture}[
  x=1mm, y=1mm, font=\small,
  cell/.style ={draw=ddpLocE, line width=0.5pt, minimum width=16mm,
                minimum height=7.6mm, inner sep=0pt, fill=ddpLoc, text=white},
  hard/.style ={-{Stealth[length=1.6mm,width=1.3mm]}, draw=ddpLocE,
                line width=0.7pt},
  soft/.style ={-{Stealth[length=1.6mm,width=1.3mm]}, draw=ddpDual,
                line width=0.5pt, densely dashed},
  stair/.style={-{Stealth[length=2.0mm,width=1.6mm]}, draw=ddpStaleE,
                line width=1.1pt, densely dashed},
  hdr/.style  ={font=\small\bfseries},
]
\def\cp{24}   % column pitch (stages)
\def\rp{14}   % row pitch    (sweeps)

% ---- headers ----
\foreach \t in {1,2,3,4}
  \node[hdr] at ({\cp*(\t-1)},{11}) {$t=\t$};
\node[hdr, anchor=south, rotate=90] at (-25,{-\rp*1.5}) {sweep $k$};

% ---- the grid of stage solves ----
\foreach \k in {1,2,3,4} {
  \node[hdr, anchor=east] at (-14,{-\rp*(\k-1)}) {$k=\k$};
  \foreach \t in {1,2,3,4}
    \node[cell] (c\k\t) at ({\cp*(\t-1)},{-\rp*(\k-1)}) {$B^{\k}[\t]$};
}

% ---- hard coupling: within a sweep, left to right ----
\foreach \k in {1,2,3,4}
  \foreach \t/\tn in {1/2, 2/3, 3/4}
    \draw[hard] (c\k\t.east) -- (c\k\tn.west);

% ---- soft coupling: from the previous sweep, one column left ----
\foreach \k/\kp in {2/1, 3/2, 4/3}
  \foreach \t/\tn in {1/2, 2/3, 3/4}
    \draw[soft] (c\kp\tn.south) -- (c\k\t.north);

% ---- the staircase: terminal information walking to stage 1 ----
\draw[stair] (c14.south) -- (c23.north);
\draw[stair] (c23.south) -- (c32.north);
\draw[stair] (c32.south) -- (c41.north);

% ---- legend: one horizontal row below the grid, so it cannot overrun the
%      two-column text width however the entries are worded ----
\def\ly{-62}
\draw[hard] (-13,{\ly}) -- (-2,{\ly});
\node[anchor=north west, align=left, text width=40mm] at (0,{\ly+3})
  {\textbf{hard}: $B[t\!-\!1]$ from \emph{this} sweep, as a constraint};
\draw[soft] (46,{\ly}) -- (57,{\ly});
\node[anchor=north west, align=left, text width=42mm] at (59,{\ly+3})
  {\textbf{soft}: $\mu[t\!+\!1]$, $B[t\!+\!1]$ \emph{frozen} from the last sweep,
   entering the objective linearly};
\draw[stair] (105,{\ly}) -- (116,{\ly});
\node[anchor=north west, align=left, text width=38mm] at (118,{\ly+3})
  {\textbf{the walk}: one terminal fact, reaching stage $1$ in sweep $T$};

\end{tikzpicture}
\caption{What the first-order scheme passes between its subproblems, drawn
for four stages over four sweeps.  Each cell is one small quadratic program.
Read one row at a time.  Within a sweep the state moves left to right along the
solid arrows and is always current --- stage $t$ is bound to the energy stage
$t\!-\!1$ has just produced, as a hard constraint.  The information travelling
the other way cannot do that.  Stage $t$ learns what the future is worth only
through $\mu[t\!+\!1]$ and $B[t\!+\!1]$ as they stood when the \emph{previous}
sweep finished, so that knowledge slides one column left per row and no faster.
The heavy staircase follows a single fact on its way through: what the terminal
stage discovers in sweep~$1$ does not reach stage~$1$ until sweep~$4$.  A
horizon of $T$ periods therefore needs about $T$ sweeps before the first stage
has heard from the last, and each sweep costs $T$ solver calls.}
\label{fig:userddp}
\end{figure*}


% Generated data: user_ddp_t3.jl.
\begin{figure*}[!t]
\centering
\begin{minipage}{0.44\textwidth}
\centering
\begin{tikzpicture}
\begin{axis}[
  width=0.88\linewidth,height=0.56\linewidth,scale only axis,
  xmin=0.7,xmax=3.3,xtick={1,2,3},xlabel={Period $t$},
  ylabel={Load factor $\lambda^t$},ymin=0.80,ymax=1.01,
  ytick={0.80,0.85,0.90,0.95,1.00},
  tick label style={font=\scriptsize},label style={font=\scriptsize},
  grid=both,major grid style={color=gridMajor},minor grid style={color=gridMinor},
  legend style={font=\scriptsize,draw=black,fill=white,fill opacity=0.85,
                text opacity=1,at={(0.02,0.98)},anchor=north west},
]
\addplot[cLoad,solid,line width=1.1pt,mark=*]
  table[col sep=comma,x=t,y=load_factor]{figures/user_ddp_t3_inputs.csv};
\addlegendentry{Load}
\end{axis}
\begin{axis}[
  width=0.88\linewidth,height=0.56\linewidth,scale only axis,
  xmin=0.7,xmax=3.3,xtick={1,2,3},axis x line=none,
  axis y line*=right,ylabel={Price [\$/kWh]},ymin=0.07,ymax=0.21,
  tick label style={font=\scriptsize},label style={font=\scriptsize},
  legend style={font=\scriptsize,draw=black,fill=white,fill opacity=0.85,
                text opacity=1,at={(0.98,0.98)},anchor=north east},
]
\addplot[cPrice,dashdotted,line width=1.1pt,mark=triangle*]
  table[col sep=comma,x=t,y=cost]{figures/user_ddp_t3_inputs.csv};
\addlegendentry{Price}
\end{axis}
\end{tikzpicture}
\end{minipage}\hfill
\begin{minipage}{0.44\textwidth}
\centering
\begin{tikzpicture}
\begin{axis}[
  width=0.88\linewidth,height=0.56\linewidth,scale only axis,
  xmin=0.5,xmax=20.5,xtick={1,4,8,12,16,20},xlabel={Completed iteration $k$},
  ylabel={Objective $J$},ymin=0.40,ymax=1.35,
  tick label style={font=\scriptsize},label style={font=\scriptsize},
  grid=both,major grid style={color=gridMajor},minor grid style={color=gridMinor},
  legend style={font=\scriptsize,draw=black,fill=white,fill opacity=0.88,
                text opacity=1,at={(0.98,0.98)},anchor=north east,
                legend columns=1,legend cell align=left},
]
\addplot[cUserDDP,solid,line width=1.0pt,mark=*,mark size=1.4pt]
  table[col sep=comma,x=iteration,y=objective]{figures/user_ddp_t3_objective.csv};
\addlegendentry{DDP}
\addplot[cFilterDDP,dashed,line width=1.1pt,mark=square*,mark size=1.5pt]
  table[col sep=comma,x=iteration,y=objective]{figures/user_ddp_t3_filter.csv};
\addlegendentry{FilterDDP}
\addplot[cCentral,densely dotted,line width=1.6pt]
  table[col sep=comma,x=iteration,y=centralized]{figures/user_ddp_t3_objective.csv};
\addlegendentry{Centralized $J^\star$}
\end{axis}
\end{tikzpicture}
\end{minipage}

\vspace{0.7em}
\begin{minipage}{0.44\textwidth}
\centering
\begin{tikzpicture}
\pgfplotsset{
  every axis/.append style={
    width=0.88\linewidth,scale only axis,
    xmin=-0.3,xmax=3.3,xtick={0,1,2,3},
    tick align=outside,
    tick label style={font=\scriptsize},
    label style={font=\scriptsize},
    grid=both,
    major grid style={color=gridMajor},
    minor grid style={color=gridMinor},
  },
}
% Battery power uses Fig. 3's physical-action convention: charging upward,
% discharging downward, with interval quantities at t-1/2.
\begin{axis}[
  name=axT3Pow,height=0.23\linewidth,
  ymin=-520,ymax=520,ytick={-450,0,450},xticklabels={},
  ylabel={$P_c\,/\,P_d$ [kW]},
  legend style={font=\scriptsize,draw=black,fill=white,fill opacity=0.88,
                text opacity=1,at={(0.02,0.98)},anchor=north west,
                row sep=-2pt,inner sep=2pt,legend cell align=left},
]
\addplot[ybar,bar width=57pt,draw=cChg!70!black,fill=cChg,fill opacity=0.75,
         line width=0.5pt,
         postaction={pattern=north west lines,pattern color=white}]
  table[col sep=comma,x=tc,y=charge_kw]{figures/user_ddp_t3_optimal_battery_power.csv};
\addlegendentry{Charging ($P_c$)}
\addplot[ybar,bar width=57pt,draw=cDis!70!black,fill=cDis,fill opacity=0.75,
         line width=0.5pt,
         postaction={pattern=north east lines,pattern color=white}]
  table[col sep=comma,x=tc,y=discharge_kw]{figures/user_ddp_t3_optimal_battery_power.csv};
\addlegendentry{Discharging ($P_d$)}
\addplot[black,line width=1.0pt,forget plot,domain=-0.3:3.3,samples=2]{0};
\end{axis}

% Stored energy follows Fig. 3: fixed B^0 in indigo, optimized B^t in purple.
\begin{axis}[
  at={(axT3Pow.below south west)},anchor=north west,yshift=-1.5mm,
  height=0.25\linewidth,
  ymin=1100,ymax=2100,ytick={1200,1500,1800,2000},
  xlabel={Time interval $t$},ylabel={$B$ [kWh]},
  legend style={font=\scriptsize,draw=black!25,fill=white,fill opacity=0.82,
                text opacity=1,at={(0.5,0.98)},anchor=north,
                legend columns=2,column sep=0.7ex,inner sep=2pt,
                legend cell align=left},
]
\addplot[ybar,bar width=57pt,draw=cB0!70!black,fill=cB0,fill opacity=0.85,
         line width=0.5pt] coordinates {(0,2000)};
\addlegendentry{Initial $B^0$}
\addplot[ybar,bar width=57pt,draw=cSOC!70!black,fill=cSOC,fill opacity=0.75,
         line width=0.5pt,restrict x to domain=1:3]
  table[col sep=comma,x=t,y=centralized_kwh]{figures/user_ddp_t3_optimal_battery.csv};
\addlegendentry{Ipopt $B^t$}
\end{axis}
\end{tikzpicture}
\end{minipage}
\caption{Three-period inputs, objective trajectories, and optimal battery actions.}
\label{fig:user_ddp_t3_results}
\end{figure*}

\clearpage

\begin{table}[!t]
\centering
\caption{SOC equality multipliers during the two-cycle}
\label{tab:user_ddp_t3_mu}
\footnotesize
\begin{tabular}{cS[table-format=1.5]S[table-format=1.5]S[table-format=1.5]}
\toprule
{$k$} & {$\lambda_{\mathrm{soc}}^{k,1}$} & {$\lambda_{\mathrm{soc}}^{k,2}$} & {$\lambda_{\mathrm{soc}}^{k,3}$} \\
\midrule
15 & 0.13500 & 0.08804 & 0.31429 \\
16 & 0.08804 & 0.23196 & 0.08804 \\
17 & 0.13500 & 0.08804 & 0.31429 \\
18 & 0.08804 & 0.23196 & 0.08804 \\
19 & 0.13500 & 0.08804 & 0.31429 \\
20 & 0.08804 & 0.23196 & 0.08804 \\
\bottomrule
\end{tabular}

\end{table}

\begin{table}[!t]
\centering
\caption{Identical-instance solver comparison}
\label{tab:user_ddp_t3_solvers}
\scriptsize
\setlength{\tabcolsep}{3pt}
\begin{tabular}{@{}lccc@{}}
\toprule
Method & Outcome & Iterations & Objective $J$ \\
\midrule
Centralized Ipopt & optimal & --- & $0.747650349$ \\
FilterDDP & optimal & $14$ & $0.747650349$ \\
First-order DDP & two-cycle & $>2000$ & $0.747650353/0.812950270$ \\
\bottomrule
\end{tabular}
\end{table}

The undamped sweep does not converge.  Through 2000 sweeps it remains in an
exact period-two cycle: the residual in \eqref{eq:app_stop} is $0.602079$, and
the objective alternates between $0.74765035$ and $0.81295027$.  The
centralized objective is $0.747650349$; on the other branch the battery dispatch
differs by as much as $0.850$~p.u. (850~kW).  Thus even this three-period
copper-plate problem exposes the instability of the first-order iteration.
On the identical data, bounds, and initialization, FilterDDP terminates normally
in 14 iterations.  Its objective differs from centralized Ipopt by
$1.48\times10^{-10}$, while its maximum battery-power and energy differences
are $6.50\times10^{-9}$~p.u. and $3.71\times10^{-9}$~p.u.h, respectively.
Consequently, the two-cycle is a property of the first-order DDP iteration, not
of an unusual or ill-posed optimization instance.
